%% -*- coding:utf-8 -*-
\chapter{Lösungen}

\section{Phrasenstrukturgrammatiken und \xbart}

\begin{enumerate}
\item Zeichnen Sie Bäume für die folgenden Phrasen. Sie dürfen das Symbol NP für Eigennamen und \nbar für
  Nomina ohne Komplemente verwenden (wie in Abbildung~\ref{Abb-das-schoene-Bild-von-Paris-HPSG}).
\eal
\ex eine Stunde vor der Ankunft des Zuges
\ex kurz nach der Ankunft in Paris
\ex das ein Lied singende Kind aus dem Allgäu
\zl
%\largerpage

\begin{figure}
\scalebox{.8}{%
\begin{forest}
sm edges
[PP
  [NP
    [Det   [eine]]
    [\nbar [Stunde]]]
  [\pbar
    [P [vor] ]
    [NP
      [Det [der]]
      [\nbar
        [N [Ankunft]]
        [NP
          [Det [des]]
          [\nbar   [Zuges]]]]]]]
\end{forest}}
\caption{Analyse von \emph{eine Stunde vor der Ankunft des Zuges}}
\end{figure}

\begin{figure}
\scalebox{.8}{%
\begin{forest}
sm edges
[PP
  [AP
    [kurz]]
  [\pbar
    [P [nach] ]
    [NP
      [Det [der]]
      [\nbar
        [\nbar [Ankunft]]
        [PP
          [\pbar
            [P [in]]
            [NP [Paris]]]]]]]]
\end{forest}}
\caption{Analyse von \emph{kurz nach der Ankunft in Paris}}
\end{figure}

\begin{figure}
\scalebox{.8}{%
\begin{forest}
sm edges
[NP
  [Det [das]]
  [\nbar
    [\nbar
      [AP
        [\abar
          [NP
            [Det [ein]]
            [\nbar [Lied]]]
          [A [singende]]]]
      [\nbar [Kind]]]
    [PP
      [\pbar
        [P [aus]]
        [NP
          [Det [dem]]
          [\nbar [Allgäu]]]]]]]
\end{forest}}
\caption{Analyse von \emph{das ein Lied singende Kind aus dem Allgäu}}
\end{figure}

\item Schreiben Sie eine Phrasenstrukturgrammatik für die Beispielsätze in (\mex{0}). Verwenden Sie
  Kategorien wie Det, A, N und P für die Phrasenstrukturregeln.

(\mex{1}) zeigt eine Grammatik, die die Bäume der vorigen Aufgabe lizenziert.
\ea
NP $\to$ Det \nbar\\
NP $\to$ \nbar PP\\
\nbar $\to$ \nbar PP\\
\nbar $\to$ AP \nbar\\
\nbar $\to$ N NP\\
%
PP $\to$ NP \pbar\\
PP $\to$ AP \pbar\\
\pbar $\to$ P NP\\
%
AP $\to$ \abar\\
\abar $\to$ NP A\\
%
\z
Sie muss mit einem Lexikon kombiniert werden, das lexikalische Einträge wie die in (\mex{1}) enthält:
\ea
Det $\to$ das\\
Det $\to$ ein\\
\nbar $\to$ Lied\\
NP $\to$ Paris\\
A $\to$ singende\\
P $\to$ aus\\
\z
\item Die Grammatik der vorigen Aufgabe kann mehr oder weniger direkt in den Computer eingegeben
  werden. Sie müssen Großbuchstaben durch Kleinbuchstaben ersetzen, Kommata zwischen den Symbolen auf
  der rechten Regelseite einfügen, \nbar durch n1 ersetzen und die lexikalischen Einträge in eckige
  Klammern einschließen:
\ea
np --> det, n1.\\
n1 --> [lied].\\
\z
\end{enumerate}
\clearpage

\section{Valenz, Argumentabfolge und Adjunktstellung}

\settowidth\jamwidth{(\ili{Deutsch})}
\begin{enumerate}
\item Geben Sie die Valenzlisten für die folgenden Wörter an:

\ea
\begin{tabular}[t]{@{}l@{~}ll@{\hspace{2em}}r@{}}
a. & laugh      & \spr \sliste{ NP[\type{nom}] }\\
b. & eat        & \spr \sliste{ NP[\type{nom}] }, \comps \sliste{ NP[\type{acc}] }\\
c. & to douse   & \spr \sliste{ NP[\type{nom}] }, \comps \sliste{ NP[\type{acc}] }\\
d. & bezichtigen& \spr \eliste, \comps \sliste{ NP[\type{nom}], NP[\type{gen}] } &(\ili{Deutsch})\\
   & accuse\\ 
e. & she & \spr \eliste, \comps \eliste\\
f. & the & \spr \eliste, \comps \eliste\\
g. & Ankunft & \spr \sliste{ Det }, \comps \sliste{ NP } & (\ili{Deutsch})\\ 
   & arrival & \spr \sliste{ Det }, \comps \eliste\\
\end{tabular}
\z
Wenn Sie sich bei der Kasuszuweisung unsicher sind, können Sie das
  Wiktionary verwenden: \url{https://de.wiktionary.org/}.

\item Zeichnen Sie Bäume für die NPn, die auch in Aufgabe~\ref{exercise-NP-PSG} auf Seite~\pageref{exercise-NP-PSG} in Kapitel~\ref{chap-psg} verwendet wurden.
\eal
\ex eine Stunde vor der Ankunft des Zuges
\ex kurz nach der Ankunft in Paris
\ex das ein Lied singende Kind aus dem Allgäu
\zl

\begin{figure}
\oneline{%
\begin{forest}
sm edges
[P\feattab{\spr   \eliste,\\
           \comps \eliste}
  [\ibox{1} N\feattab{\spr   \eliste,\\
                      \comps \eliste}
    [\ibox{2} Det   [eine]]
    [N\feattab{\spr   \sliste{ \ibox{2} },\\
               \comps \eliste} [Stunde]]]
  [P\feattab{\spr   \sliste{ \ibox{1} },\\
             \comps \eliste}
    [P\feattab{\spr   \sliste{ \ibox{1} },\\
               \comps \sliste{ \ibox{3} }} [vor] ]
    [\ibox{3} N\feattab{\spr   \eliste,\\
                         \comps \eliste}
      [\ibox{4} Det [der]]
      [N\feattab{\spr   \sliste{ \ibox{4} },\\
                 \comps \eliste}
        [N\feattab{\spr   \sliste{ \ibox{4} },\\
                   \comps \sliste{ \ibox{5} }} [Ankunft]]
        [\ibox{5} N\feattab{\spr   \eliste,\\
                            \comps \eliste}
          [\ibox{6} Det [des]]
          [N\feattab{\spr   \sliste{ \ibox{6} },\\
                     \comps \eliste}   [Zuges]]]]]]]
\end{forest}}
\caption{Analyse von \emph{eine Stunde vor der Ankunft des Zuges}}
\end{figure}

\begin{figure}
\begin{forest}
sm edges
[P\feattab{\spr   \eliste,\\
           \comps \eliste}
  [\ibox{1} A\feattab{\spr   \eliste,\\
                      \comps \eliste}
    [kurz]]
  [P\feattab{\spr   \sliste{ \ibox{1} },\\
             \comps \eliste}
    [P\feattab{\spr   \sliste{ \ibox{1} },\\
               \comps \sliste{ \ibox{2} }} [nach] ]
    [\ibox{2} N\feattab{\spr   \eliste,\\
                        \comps \eliste}
      [\ibox{3} Det [der]]
      [N\feattab{\spr   \sliste{ \ibox{3} },\\
                 \comps \eliste}
        [\ibox{4} N\feattab{\spr   \sliste{ \ibox{3} },\\
                            \comps \eliste} [Ankunft]]
        [P\feattab{\textsc{mod} \ibox{4},\\
                   \spr   \eliste,\\
                   \comps \eliste}
          [P\feattab{\textsc{mod} \ibox{4},\\
                     \spr   \eliste,\\
                     \comps \sliste{ \ibox{5} }} [in]]
          [\ibox{5} N\feattab{\spr   \eliste,\\
                     \comps \eliste} [Paris]]]]]]]
\end{forest}
\caption{Analyse von \emph{kurz nach der Ankunft in Paris}}
\end{figure}

\begin{figure}
%\oneline{%
\scalebox{.85}{%
\begin{forest}
sm edges
[N\feattab{\spr   \eliste,\\
           \comps \eliste}
  [\ibox{1} Det [das]]
  [N\feattab{\spr   \sliste{ \ibox{1} },\\
             \comps \eliste}
    [\ibox{2} N\feattab{\spr   \sliste{ \ibox{1} },\\
               \comps \eliste}
      [A\feattab{\spr   \eliste,\\
                 \comps \eliste}, tier=3
        [\ibox{4} N\feattab{\spr   \eliste,\\
                   \comps \eliste}
          [\ibox{6} Det [ein]]
          [N\feattab{\spr   \sliste{ \ibox{6} },\\
                     \comps \eliste} [Lied]]]
        [A\feattab{\spr   \eliste{ },\\
                   \comps \sliste{ \ibox{4} }} [singende]]]
      [N\feattab{\spr   \sliste{ \ibox{1} },\\
                 \comps \eliste}, tier=3 [Kind]]]
    [P\feattab{\textsc{mod} \ibox{2},\\
               \spr   \eliste,\\
               \comps \eliste}
      [P\feattab{\textsc{mod} \ibox{2},\\
                 \spr   \eliste,\\
                 \comps \sliste{ \ibox{3} }},tier=3 [aus]]
      [\ibox{3} N\feattab{\spr   \eliste,\\
                          \comps \eliste},tier=3
        [\ibox{5} Det [dem]]
        [N\feattab{\spr   \sliste{ \ibox{5} },\\
                   \comps \eliste} [Allgäu]]]]]]]
\end{forest}}
\caption{Analyse von \emph{das ein Lied singende Kind aus dem Allgäu}}
\end{figure}


\item Zeichnen Sie Bäume für die folgenden Beispiele. NPn können abgekürzt werden.

\eal
\ex weil Aicke dem Kind ein Buch schenkt
\ex weil dem Kind solch ein Buch niemand schenkt
\ex because Kim gave a book to him
\ex Sandy saw this yesterday.
\ex
\gll at Bjarne læste bogen\\
     dass Bjarne las Buch.\textsc{def}\\\danish
\glt `dass Bjarne das Buch las'
\zl


Die Bäume mit den Lösungen sind im
Folgenden angegeben. Abbildung~\ref{fig-weil-dem-Kind-solch-ein-buch-niemand-schenkt} unterscheidet sich von
Abbildung~\ref{fig-weil-Aicke-dem-Kind-ein-buch-schenkt} in der Art, wie die Elemente in der \compsl
nummeriert sind, aber in jedem Fall ist die Reihenfolge der Elemente in der \compsl von \emph{schenkt}
\sliste{ \npnom, \npdat, \npacc }. Die unterschiedliche Nummerierung ist auf die Reihenfolge zurückzuführen, in der
die Elemente kombiniert werden. Wenn die Nummerierung durchgängig von oben nach unten vorgenommen wird, ergibt sich
Abbildung~\ref{fig-weil-dem-Kind-solch-ein-buch-niemand-schenkt}. Wenn man bei der Vergabe der Nummern großzügiger
ist, kann dieselbe Situation wie in
Abbildung~\ref{fig-weil-dem-Kind-solch-ein-buch-niemand-schenkt-different-numbering} dargestellt werden. Abbildung~\ref{fig-weil-dem-Kind-solch-ein-buch-niemand-schenkt-different-numbering}
hat dieselbe Nummerierung in der Valenzliste wie Abbildung~\ref{fig-weil-Aicke-dem-Kind-ein-buch-schenkt}
und ist deshalb vielleicht leichter zu erfassen.

\begin{figure}
%\centerfit{%
\scalebox{.85}{%
\begin{forest}
sm edges
[{C[\spr \eliste, \comps \eliste]}
  [{C[\spr \eliste, \comps \sliste{ \ibox{1} }]} [weil]]
  [{\ibox{1} V[\spr \eliste, \comps \eliste]}
      [{\ibox{2} NP[\type{nom}]} [Aicke] ]
      [V\feattab{
            \spr \eliste, \comps \sliste{ \ibox{2} }}
        [{\ibox{3} NP[\type{dat}]} [dem Kind,roof] ]
        [V\feattab{
              \spr \eliste, \comps \sliste{ \ibox{2}, \ibox{3} }}
          [{\ibox{4} NP[\type{acc}]} [ein Buch,roof] ]
          [V\feattab{
              \spr \sliste{  },\\
              \comps \sliste{ \ibox{2}, \ibox{3}, \ibox{4} }} [schenkt] ]]
] ] ]
\end{forest}}
\caption{\label{fig-weil-Aicke-dem-Kind-ein-buch-schenkt}Die Analyse von \emph{weil Aicke dem Kind
    ein Buch schenkt}}
\end{figure}

\begin{figure}
\centerfit{%
\begin{forest}
sm edges
[{C[\spr \eliste, \comps \eliste]}
  [{C[\spr \eliste, \comps \sliste{ \ibox{1} }]} [weil]]
  [{\ibox{1} V[\spr \eliste, \comps \eliste]}
      [{\ibox{2} NP[\type{acc}]} [solch ein Buch,roof] ]
      [V\feattab{
            \spr \eliste, \comps \sliste{ \ibox{2} }}
        [{\ibox{3} NP[\type{dat}]} [dem Kind,roof] ]
        [V\feattab{
              \spr \eliste, \comps \sliste{ \ibox{3}, \ibox{2} }}
          [{\ibox{4} NP[\type{nom}]} [niemand] ]
          [V\feattab{
              \spr \sliste{  },\\
              \comps \sliste{ \ibox{4}, \ibox{3}, \ibox{2} }} [schenkt] ]]
] ] ]
\end{forest}}
\caption{\label{fig-weil-dem-Kind-solch-ein-buch-niemand-schenkt}Die Analyse von \emph{weil dem Kind
    solch ein Buch niemand schenkt}}
\end{figure}

\begin{figure}
\centerfit{%
\begin{forest}
sm edges
[{C[\spr \eliste, \comps \eliste]}
  [{C[\spr \eliste, \comps \sliste{ \ibox{1} }]} [weil]]
  [{\ibox{1} V[\spr \eliste, \comps \eliste]}
      [{\ibox{4} NP[\type{acc}]} [solch ein Buch,roof] ]
      [V\feattab{
            \spr \eliste, \comps \sliste{ \ibox{4} }}
        [{\ibox{3} NP[\type{dat}]} [dem Kind,roof] ]
        [V\feattab{
              \spr \eliste, \comps \sliste{ \ibox{3}, \ibox{4} }}
          [{\ibox{2} NP[\type{nom}]} [niemand] ]
          [V\feattab{
              \spr \sliste{  },\\
              \comps \sliste{ \ibox{2}, \ibox{3}, \ibox{4} }} [schenkt] ]]
] ] ]
\end{forest}}
\caption{\label{fig-weil-dem-Kind-solch-ein-buch-niemand-schenkt-different-numbering}Die Analyse von \emph{weil dem Kind
    solch ein Buch niemand schenkt}}
\end{figure}


\begin{figure}
\centerfit{%
\begin{forest}
sm edges
[{C[\spr \eliste, \comps \eliste]}
  [{C[\spr \eliste, \comps \sliste{ \ibox{1} }]} [because]]
  [{\ibox{1} V[\spr \eliste, \comps \eliste]}
      [{\ibox{2} NP[\type{nom}]} [Kim] ]
      [V\feattab{
            \spr \sliste{ \ibox{2} }, \comps \eliste}
        [V\feattab{
                \spr \sliste{ \ibox{2} }, \comps \sliste{ \ibox{3} }}
           [V\feattab{
              \spr \sliste{ \ibox{2} },\\
              \comps \sliste{ \ibox{4}, \ibox{3} }} [gave] ]
            [{\ibox{4} NP[\type{acc}]} [a book, roof] ] ] 
        [{\ibox{3} PP[\type{to}]} [to him,roof] ] ] ] ]
\end{forest}}
\caption{\label{fig-because-Kim-gave-the-book-to-him}Die Analyse von \emph{because Kim gave a book to him}}
\end{figure}

\begin{figure}
\centerfit{%
\begin{forest}
sm edges
[{V[\spr \eliste, \comps \eliste]}
   [{\ibox{1} NP[\type{nom}]} [Sandy] ]
   [V\feattab{
      \spr \sliste{ \ibox{1} }, \comps \eliste}
     [\ibox{2} V\feattab{
        \spr \sliste{ \ibox{1} }, \comps \eliste}
       [V\feattab{
           \spr \sliste{ \ibox{1} },\\
           \comps \sliste{ \ibox{3} }} [saw] ]
          [{\ibox{3} NP[\type{acc}]} [this] ] ]
     [{Adv[\textsc{mod} \ibox{2} VP]} [yesterday] ] ] ]
\end{forest}}
\caption{\label{fig-Sandy-saw-this-yesterday}Analyse von \emph{Sandy saw this yesterday.}}
\end{figure}

\begin{figure}
%\centerfit{%
\scalebox{.85}{%
\begin{forest}
sm edges
[{C[\spr \eliste, \comps \eliste]}
  [{C[\spr \eliste, \comps \sliste{ \ibox{1} }]} [at;dass]]
  [{\ibox{1} V[\spr \eliste, \comps \eliste]}
    [{\ibox{2} NP[\type{nom}]} [Bjarne;Bjarne] ]
    [V\feattab{
       \spr \sliste{ \ibox{2} }, \comps \eliste}
       [V\feattab{
           \spr \sliste{ \ibox{2} },\\
           \comps \sliste{ \ibox{3} }} [læste;las] ]
          [{\ibox{3} NP[\type{acc}]} [bogen;Buch.\textsc{def}] ] ] ] ]
\end{forest}}
\caption{\label{fig-at-bjarne-laeste-bogen}Analyse von \emph{at Bjarne læste bogen} `dass Bjarne
  das Buch las'}
\end{figure}



\end{enumerate}

%\clearpage

\section{Der Verbalkomplex}

\begin{enumerate}
\item Skizzieren Sie die Analyse der Verbalkomplexe in den folgenden Beispielen:
\eal
\ex dass sie darüber lachen muss
\ex dass sie darüber hat lachen müssen
\ex dass sie darüber wird haben lachen müssen
\zl
Sie dürfen die \spr-Werte weglassen, da sie für alle deutschen\il{Deutsch} Verben ohnehin die leere Liste sind.


\begin{figure}
\scalebox{.9}{%
\begin{forest}
sm edges
[C\feattab{%\spr \eliste,\\
             \comps \eliste}
  [C\feattab{%\spr \eliste,\\
             \comps \sliste{ \ibox{1} }} [dass]]
  [{\ibox{1} V\feattab{%\spr \eliste,\\
                       \comps \eliste}}
     [{\ibox{2} NP[nom]} [sie]]
     [V\feattab{%\spr  \eliste,\\
                \comps \sliste{ \ibox{2} }}
       [\ibox{3} PP [darüber]]
       [V\feattab{%\spr \eliste,\\
                  \comps \sliste{ \ibox{2}, \ibox{3} }}
         [\ibox{4} V\feattab{\subj \sliste{ \ibox{2} },\\
                             %\spr \eliste,\\
                             \comps \sliste{ \ibox{3} }} [lachen]]
         [V\feattab{%\spr \eliste,\\
                    \comps \sliste{ \ibox{2}, \ibox{3}, \ibox{4} }} [muss]]]
]]]
\end{forest}}
\caption{Analyse von \emph{dass sie darüber lachen muss}}
\end{figure}


\begin{figure}
%\oneline{%
\scalebox{.9}{%
\begin{forest}
sm edges
[C\feattab{%\spr \eliste,\\
             \comps \eliste}
  [C\feattab{%\spr \eliste,\\
             \comps \sliste{ \ibox{1} }} [dass]]
  [{\ibox{1} V\feattab{%\spr \eliste,\\
                       \comps \eliste}}
     [{\ibox{2} NP[nom]} [sie]]
     [V\feattab{%\spr  \eliste,\\
                 \comps \sliste{ \ibox{2} }}
       [\ibox{3} PP [darüber]]
       [V\feattab{%\spr \eliste,\\
                  \comps \sliste{ \ibox{2}, \ibox{3} }}
         [V\feattab{%\spr \eliste,\\
                    \comps \sliste{ \ibox{2}, \ibox{3}, \ibox{4} }} [hat]]
           [\ibox{4} V\feattab{\subj \sliste{ \ibox{2} },\\
                               %\spr \eliste,\\
                               \comps \sliste{ \ibox{3} }}
             [\ibox{5} V\feattab{\subj \sliste{ \ibox{2} },\\
                                 %\spr \eliste,\\
                                 \comps \sliste{ \ibox{3} }} [lachen]]
             [V\feattab{\subj \sliste{ \ibox{2} },\\
                        %\spr \eliste,\\
                        \comps \sliste{ \ibox{3}, \ibox{5} }} [müssen]]]]]]]
\end{forest}}
\caption{Analyse von \emph{dass sie darüber hat lachen müssen}}
\end{figure}


\begin{figure}
\oneline{%
\begin{forest}
sm edges
[C\feattab{%\spr \eliste,\\
             \comps \eliste}
  [C\feattab{%\spr \eliste,\\
             \comps \sliste{ \ibox{1} }} [dass]]
  [{\ibox{1} V\feattab{%\spr \eliste,\\
                       \comps \eliste}}
     [{\ibox{2} NP[nom]} [sie]]
     [V\feattab{%\spr  \eliste,\\
                 \comps \sliste{ \ibox{2} }}
       [\ibox{3} PP [darüber]]
       [V\feattab{%\spr \eliste,\\
                  \comps \sliste{ \ibox{2}, \ibox{3} }}
         [V\feattab{%\spr \eliste,\\
                    \comps \sliste{ \ibox{2}, \ibox{3}, \ibox{4} }} [wird]]
           [\ibox{4} V\feattab{\subj \sliste{ \ibox{2} },\\
                              %\spr \eliste,\\
                               \comps \sliste{ \ibox{3} }}
             [V\feattab{\subj \sliste{ \ibox{2} },\\
                       %\spr \eliste,\\
                        \comps \sliste{ \ibox{3}, \ibox{5} }} [haben]]
             [\ibox{5} V\feattab{\subj \sliste{ \ibox{2} },\\
                                 %\spr \eliste,\\
                                 \comps \sliste{ \ibox{3} }}
               [\ibox{6} V\feattab{\subj \sliste{ \ibox{2} },\\
                                   %\spr \eliste,\\
                                   \comps \sliste{ \ibox{3} }} [lachen]]
               [V\feattab{\subj \sliste{ \ibox{2} },\\
                         %\spr \eliste,\\
                          \comps \sliste{ \ibox{3}, \ibox{6} }} [müssen]]]]]]]]
\end{forest}}
\caption{Analyse von \emph{dass sie darüber wird haben lachen müssen}}
\end{figure}



\end{enumerate}
\clearpage


\section{Verbstellung: Verberst- und Verbzweitstellung}



\begin{enumerate}
\item Klassifizieren Sie die germanischen\il{Germanisch} Sprachen nach ihrer grundlegenden Konstituentenabfolge (SVO, SOV, VSO,
  \ldots) und nach V2 unter der Annahme, dass Sie wissen, dass eines der folgenden Muster in der Sprache vorliegt:
\eal
\label{ex-v2-task-solution}
\ex 
\label{ex-acc-aux-nom-v-dat}
NP[acc] V-Aux NP[nom] V NP[dat]   \hfill  V2 SVO 
\ex
\label{ex-acc-aux-nom-dat-v} 
NP[acc] V-Aux NP[nom] NP[dat] V   \hfill  V2 SOV
\ex 
\label{ex-acc-nom-v-acc}
NP[acc] NP[nom] V NP[acc]         \hfill $-$V2 SVO
\ex 
\label{ex-acc-nom-aux-v-acc}
NP[acc] NP[nom] V-Aux V NP[acc]   \hfill $-$V2 SVO
\ex 
\label{ex-acc-aux-nom-v-pp}
NP[acc] V-Aux NP[nom] V PP        \hfill nicht klassifizierbar
\zl

%\largerpage[1]
\noindent
Das Muster in (\ref{ex-acc-aux-nom-v-dat}) kann nicht das \ili{Englisch}e sein, da das \ili{Englisch}e keinen Dativ hat. Folglich handelt es sich um eine V2"=Sprache.
Das Dativobjekt folgt dem Verb, also muss es eine SVO"=Sprache sein. Ein Beispiel wäre das \ili{Isländisch}e:
\ea
\gll Bókina          hafa ég       gefið honum.\\
     Buch.das.\ACC{} haben ich.\NOM{} gegeben er.\DAT\\\icelandic
\glt `Ich habe ihm das Buch gegeben.'
\z

(\ref{ex-acc-aux-nom-dat-v}) hat ein Hilfsverb und zwei NPn, gefolgt von einem Verb. Da das Dativobjekt in einer
SVO"=Sprache dem Verb folgen würde, muss es eine SOV"=Sprache sein. Da alle germanischen\il{Germanisch} SOV"=Sprachen auch V2"=Sprachen
sind, muss (\ref{ex-acc-aux-nom-dat-v}) eine V2"=Sprache sein. Das \ili{Deutsch}e und das \ili{Niederländisch}e wären Beispiele.

\ea Den Roman hat jemand dem Kind gegeben.
\z

Sieht man von Mehrfachvorfeldbesetzungen im \ili{Deutsch}en ab \citep{Mueller2003b}, so muss (\ref{ex-acc-nom-v-acc}) ein Nicht-V2"=Muster sein. Die Sprache kann nur das
\ili{Englisch}e sein:
\ea
This book, Kim gave Sandy.
\z
Aus demselben Grund ist (\ref{ex-acc-nom-aux-v-acc}) nicht-V2 und SVO. Die Sprache muss das \ili{Englisch}e sein:

\ea
This book, Kim had given Sandy.
\z

\noindent
Das Muster in (\ref{ex-acc-aux-nom-v-pp}) kann hinsichtlich V2 und
SOV/SVO nicht eindeutig klassifiziert werden. Da PPn leicht extraponiert werden können, könnte es eine SOV"=\pagebreak{}Sprache mit Extraposition sein (\zb{}
das \ili{Deutsch}e), oder es könnte das \ili{Englisch}e mit Fragebildung (Residual-V2) sein:

\eal
\ex Wen hat Aicke gesehen bei der Demonstration?
\ex Who did Kim see during the rally?
\zl

%\largerpage
\item Skizzieren Sie die Analyse für die folgenden Beispiele. Verwenden Sie die in diesem Kapitel verwendeten
  Abkürzungen; das heißt, gehen Sie nicht auf die Details bezüglich \spr und \compsw ein, sondern verwenden Sie S, VP und V$'$.
\eal
\ex
\gll Arbejder Bjarne ihærdigt  på bogen?\\
     arbeitet Bjarne ernsthaft an Buch.\textsc{def}\\\danish
\glt `Arbeitet Bjarne ernsthaft an dem Buch?'
\ex Arbeitet Bjarne ernsthaft an dem Buch?
\ex Wird sie darüber nachdenken?
\zl

\begin{figure}
\scalebox{.95}{%
\begin{forest}
sm edges
[S
  [V\sliste{ S//V }
    [V [arbejder;arbeitet]]]
  [S//V
     [NP [Bjarne;Bjarne]]
     [VP//V
       [Adv [ihærdigt;ernsthaft]]
       [VP//V
         [V//V [\trace]]
         [PP [på bogen;{an Buch.\textsc{def}},roof]]]]]]
\end{forest}}
\caption{Analyse von \emph{Arbejder Bjarne ihærdigt på bogen?} `Arbeitet Bjarne ernsthaft an dem Buch?'}
\end{figure}

\begin{figure}
\begin{forest}
sm edges
[S
  [V\sliste{ S//V }
    [V [arbeitet]]]
  [S//V
     [NP [Bjarne]]
     [\vbarDSLv
       [Adv [ernsthaft]]
       [\vbarDSLv
         [PP [an dem Buch,roof]]
         [V//V [\trace]]]]]]
\end{forest}
\caption{Analyse von \emph{Arbejder Bjarne ihærdigt på bogen?} `Arbeitet Bjarne ernsthaft an dem Buch?'}
\end{figure}

\begin{figure}
\begin{forest}
sm edges
[S
  [V\sliste{ S//V }
    [V
      [wird]]]
  [S//V
     [NP [sie]]
     [\vbarDSLv
       [PP [darüber]]
       [V//V
         [V [nachdenken]]
         [V//V [\trace]]]
]]]
\end{forest}
\caption{Analyse von \emph{Wird sie darüber nachdenken?}}
\end{figure}

\clearpage

\item Skizzieren Sie die Analyse für die folgenden Beispiele. Verwenden Sie die Valenzmerkmale \spr und \comps
  anstelle der Abkürzungen S, VP und V$'$. Da der Wert von \spr im \ili{Deutsch}en immer die
  leere Liste ist, dürfen Sie ihn in den deutschen\il{Deutsch} Beispielen weglassen. NPn und PPn können als NP bzw. PP abgekürzt werden.
\eal
\ex dass sie darüber nachdenkt
\ex dass sie darüber nachdenken wird
\ex Wird sie darüber nachdenken?
\zl

\eal
\ex
\gll Arbejder Bjarne ihærdigt  på bogen?\\
     arbeitet Bjarne ernsthaft an Buch.\textsc{def}\\\jambox{(\ili{Dänisch})}
\glt `Arbeitet Bjarne ernsthaft an dem Buch?'
\ex Arbeitet Bjarne ernsthaft an dem Buch?
\zl


\begin{figure}
\begin{forest}
sm edges
[C\feattab{%\spr \eliste,\\
             \comps \eliste}
  [C\feattab{%\spr \eliste,\\
             \comps \sliste{ \ibox{1} }} [dass]]
  [{\ibox{1} V\feattab{%\spr \eliste,\\
                       \comps \eliste}}
     [{\ibox{2} NP[nom]} [sie]]
     [V\feattab{%\spr  \eliste,\\
                 \comps \sliste{ \ibox{2} }}
       [\ibox{3} PP [darüber]]
       [V\feattab{%\spr \eliste,\\
                  \comps \sliste{ \ibox{2}, \ibox{3} }} [nachdenkt]]]]]
\end{forest}
\caption{Analyse von \emph{dass sie darüber nachdenkt}}
\end{figure}


\begin{figure}
\begin{forest}
sm edges
[C\feattab{%\spr \eliste,\\
             \comps \eliste}
  [C\feattab{%\spr \eliste,\\
             \comps \sliste{ \ibox{1} }} [dass]]
  [{\ibox{1} V\feattab{%\spr \eliste,\\
                       \comps \eliste}}
     [{\ibox{2} NP[nom]} [sie]]
     [V\feattab{%\spr  \eliste,\\
                 \comps \sliste{ \ibox{2} }}
       [\ibox{3} PP [darüber]]
       [V\feattab{%\spr \eliste,\\
                  \comps \sliste{ \ibox{2}, \ibox{3} }}
         [\ibox{4} V\feattab{\subj \sliste{ \ibox{2} },\\
                             %\spr \eliste,\\
                             \comps \sliste{ \ibox{3} }} [nachdenken]]
         [V\feattab{%\spr \eliste,\\
                    \comps \sliste{ \ibox{2}, \ibox{3}, \ibox{4} }} [wird]]]
]]]
\end{forest}
\caption{Analyse von \emph{dass sie darüber nachdenken wird}}
\end{figure}


\begin{figure}
\begin{forest}
sm edges
[V\feattab{%\spr \eliste,\\
           \comps \eliste}
  [V\feattab{%\spr \eliste,\\
             \comps \sliste{ \ibox{1} }}
    [V
      [wird]]]
  [{\ibox{1} V//V\feattab{%\spr \eliste,\\
                       \comps \eliste}}
     [{\ibox{2} NP[nom]} [sie]]
     [V//V\feattab{%\spr  \eliste,\\
                 \comps \sliste{ \ibox{2} }}
       [\ibox{3} PP [darüber]]
       [V//V\feattab{%\spr \eliste,\\
                  \comps \sliste{ \ibox{2}, \ibox{3} }}
         [\ibox{4} V\feattab{\subj \sliste{ \ibox{2} },\\
                             %\spr \eliste,\\
                             \comps \sliste{ \ibox{3} }} [nachdenken]]
         [V//V\feattab{%\spr \eliste,\\
                    \comps \sliste{ \ibox{2}, \ibox{3}, \ibox{4} }} [\trace]]]
]]]
\end{forest}
\caption{Analyse von \emph{Wird sie darüber nachdenken?}}
\end{figure}



%\gll Arbejder Bjarne ihærdigt  på bogen?\\
%     works    Bjarne seriously at book.\textsc{def}\\\jambox{(\ili{Danish})}


\begin{figure}
\begin{forest}
sm edges
[V\feattab{\comps \eliste}
  [V\feattab{\comps \sliste{ \ibox{1} }}
    [V [arbejder;arbeitet]]]
  [\ibox{1} V//V\feattab{\spr   \eliste,\\
                         \comps \eliste}
     [\ibox{2} NP [Bjarne;Bjarne]]
     [V//V\feattab{\spr   \sliste{ \ibox{2} },\\
                    \comps \eliste}
       [{Adv[\textsc{mod} \ibox{3} VP]} [ihærdigt;ernsthaft]]
       [\ibox{3} V//V\feattab{\spr   \sliste{ \ibox{2} },\\
                      \comps \eliste}
         [V//V\feattab{\spr \sliste{ \ibox{2} },\\
                       \comps \sliste{ \ibox{4} }} [\trace]]
         [\ibox{4} PP [på bogen;{an Buch.\textsc{def}},roof]]]]]]
\end{forest}
\caption{Analyse von \emph{Arbejder Bjarne ihærdigt  på bogen?} `Arbeitet Bjarne ernsthaft an dem Buch?'}
\end{figure}

\begin{figure}
\begin{forest}
sm edges
[V\feattab{\comps \eliste}
  [V\feattab{\comps \sliste{ \ibox{1} }}
    [V [arbeitet]]]
  [\ibox{1} V//V\feattab{%\spr   \eliste,\\
                         \comps \eliste}
     [\ibox{2} NP [Bjarne]]
     [V//V\feattab{%\spr   \eliste,\\
                    \comps \sliste{ \ibox{2} }}
       [{Adv[\textsc{mod} \ibox{3} V[\textsc{ini}$-$]]} [ernsthaft]]
       [\ibox{3} V//V\feattab{%\spr   \sliste{  },\\
                      \comps \sliste{ \ibox{2} }}
         [\ibox{4} PP [an dem Buch,roof]]
         [V//V\feattab{%\spr \sliste{  },\\
                       \comps \sliste{ \ibox{2}, \ibox{4} }} [\trace]]]]]]
\end{forest}
\caption{Analyse von \emph{Arbeitet Bjarne ernsthaft an dem Buch?}}
\end{figure}


\item Skizzieren Sie die Analyse der folgenden Beispiele. NPn dürfen abgekürzt werden. Valenzmerkmale sollen
  nicht angegeben werden, sondern stattdessen Knotenlabels wie V, V$'$, VP und S. Falls nichtlokale
  Abhängigkeiten beteiligt sind, kennzeichnen Sie diese mit dem Symbol `/'.
\eal
\ex Such books, I like.
\ex Solche Bücher mag ich.
\ex
\gll Boger som det elsker jeg.\\
     Bücher wie diese mag ich\\\danish
\glt `Ich mag solche Bücher.'
\zl

\begin{figure}
\scalebox{1}{%
\begin{forest}
sm edges
[S
  [NP$_i$ [such books,roof] ]
  [S/NP 
    [NP [I] ] 
    [VP/NP  
      [V [like] ]
      [NP/NP [\trace$_i$] ] ] ] ]
\end{forest}}
\caption{Analyse von \emph{Such books, I like.}}
\end{figure}

\begin{figure}
\scalebox{1}{%
\begin{forest}
sm edges
[S
  [NP$_i$ [solche Bücher, roof] ]
  [S/NP
     [V \sliste{ S//V }
        [V [mag$_j$] ] ]
     [S//\vSLASHnp
        [NP/NP [\trace$_i$] ]
        [\vbarDSLv
           [NP [ich] ]
           [V//V [\_$_j$] ] ] ] ] ] ]
\end{forest}}
\caption{Analyse des deutschen Beispiels \emph{Solche Bücher mag ich.}}
\end{figure}


\begin{figure}
\begin{forest}
sm edges
[S
   [NP$_i$ [boger som det;Bücher wie diese,roof] ]
      [S/NP
         [V \sliste{ S//V }
           [V [elsker$_j$;mag] ] ]
           [S//\vSLASHnp
             [NP [jeg;ich] ]
             [VP//\vSLASHnp
               [V//V  [\_$_j$] ]
               [NP/NP [\trace$_i$ ] ] ] ] ] ] ] 
\end{forest}
\caption{\label{fig-boger-som-et-elsker-jeg}Analyse des dänischen Beispiels \emph{Boger som det elsker jeg.} `Ich mag solche Bücher.'}
\end{figure}

\item Skizzieren Sie die Analyse des Beispiels in (\mex{0}b) und nehmen Sie das Merkmal \comps mit auf.

\begin{figure}
\scalebox{1}{%
\begin{forest}
sm edges
[S
  [NP$_i$ [solche Bücher, roof] ]
  [S/NP
     [{V[\comps \sliste{ \ibox{1} }]}
        [V [mag$_j$] ] ]
     [{\ibox{1} S//\vSLASHnp[\comps \sliste{ }]}
        [\ibox{3} NP/NP [\trace$_i$] ]
        [{\vbarDSLv[\comps \sliste{ \ibox{3} }]}
           [\ibox{2} NP [ich] ]
           [{V//V[\comps \sliste{ \ibox{2}, \ibox{3} }]}  [\_$_j$] ] ] ] ] ] ]
\end{forest}}
\caption{Analyse des deutschen Beispiels \emph{Solche Bücher mag ich.}}
\end{figure}


\end{enumerate}

\clearpage

\section{Passiv}

\begin{enumerate}
\item Welche NPn in (\mex{1}) haben strukturellen und welche lexikalischen Kasus?
\eal
\ex Der Junge lacht.
\ex Mich friert.
\ex Er zerstört die Umwelt.
\ex Das dauert ein ganzes Jahr.
\ex Er hat nur einen Tag dafür gebraucht.
\ex Er denkt an den morgigen Tag.
\zl

Alle Nominative in (\mex{0}) sind strukturelle Kasus. Die Akkusative in (\mex{0}b, d, f) sind lexikalisch,
die in (\mex{0}c, e) sind strukturell.

\item Geben Sie \argst-Listen für die folgenden Verben an. Geben Sie die \argstl mit der maximalen Anzahl von Argumenten an.
\eal
\ex show, eat, meet \english
\ex zeigen, essen, begegnen, treffen \german
%\ex vise `show', spise `eat', møde `meet' \danish
%\ex sýna `show', eta `eat', mæta `meet', hitta `meet' \icelandic
\zl

\eal
\ex \emph{show}: \sliste{ NP[\str], NP[\str], NP[\type{lacc}] }
\ex \emph{eat}: \sliste{ NP[\str], NP[\str] }
\ex \emph{meet}: \sliste{ NP[\str], NP[\str] }
\zl
\eal
\ex \emph{zeigen}: \sliste{ NP[\str], NP[\type{ldat}], NP[\str] }
\ex \emph{essen}: \sliste{ NP[\str], NP[\str] }
\ex \emph{begegnen}: \sliste{ NP[\str], NP[\type{ldat}] }
\ex \emph{treffen}: \sliste{ NP[\str], NP[\str] }
\zl
Wenn Sie sich beim Kasus unsicher sind, können Sie das
  Wiktionary verwenden: \url{https://de.wiktionary.org/}.

\item Zeichnen Sie den Analysebaum für den folgenden Satz:
\ea
that the box was opened
\z
Geben Sie bitte Valenzmerkmale (\spr und \comps) und Wortartinformationen an. Sie dürfen
die NP mit einem Dreieck abkürzen.


\begin{figure}
\begin{forest}
sm edges
[{C[\comps \eliste]}
  [{C[\comps \sliste{ \ibox{1} }]} [that]]
  [{\ibox{1} V\feattab{\spr \eliste,\\
                       \comps \eliste}}
     [{\ibox{2} NP[nom]} [the box, roof]]
     [V\feattab{\spr  \sliste{ \ibox{2} },\\
                 \comps \eliste}
       [V\feattab{\spr  \sliste{ \ibox{2} },\\
                  \comps \sliste{ \ibox{3} }} [was]]
       [{\ibox{3} V\feattab{\spr \sliste{ \ibox{2} },\\
                            \comps \eliste}} [opened]]]]]
\end{forest}
\caption{Analyse des Passivsatzes \emph{that the box was opened}}
\end{figure}
Das transitive Verb \emph{open} nimmt ein Subjekt und ein Objekt. Die \argst-Liste enthält zwei NPn mit
strukturellem Kasus. Die Passiv-Lexikonregel entfernt ein Argument. Für das Passivpartizip
bleibt uns damit ein Element auf der \argstl. Dieses Element wird auf die \sprl von
\emph{opened} abgebildet. Das Passivhilfsverb nimmt eine VP in Passivform und übernimmt deren Element aus
\spr. Nach der Kombination von Hilfsverb und Passiv-VP haben wir die VP \emph{was opened}, die immer noch
einen Spezifikator selegiert. Die NP \emph{the box} fungiert als Spezifikator, und die Kombination von
\emph{the box} und \emph{was opened} ist ein vollständiger Satz.

\end{enumerate}

%% \begin{figure}
%% \begin{forest}
%% sm edges
%% [CP
%%   [C [dass;that]]
%%   [V
%%      [\ibox{1} {NP[nom]} [der Kasten;the box, roof]]
%%      [V
%%        [{\ibox{2} V[\comps \ibox{3} \sliste{ \ibox{1} }]} [geöffnet;opened]]
%%        [{V[\comps \ibox{3} $\oplus$ \sliste{ \ibox{2} }] [wurde;was]]]]]
%% \end{forest}
%% \end{figure}

\section{Satztypen und Expletiva}

\begin{enumerate}
\item Analysieren Sie die Interrogativsätze in (\mex{1}):
\eal
\ex Ich weiß, wen Kim kennt.
\ex
\gll Jeg ved, hvem det kende Kim\\
     ich weiß wer  \expl{} kennt Kim\\\danish
\glt `Ich weiß, wer Kim kennt.'
\zl


\begin{figure}
\centerline{\begin{forest}
sm edges
[S
       [{NP[\sacc]} [wen] ]
       [{S/NP[\sacc]}
         [{NP[\sacc]/NP[\sacc]}  [\trace] ]
         [V$'$
           [{NP[\snom]} [Kim] ]
           [V [kennt] ] ] ] ]
\end{forest}}
\caption{Analyse von \emph{wen Kim kennt}}
\end{figure}


\begin{figure}
\centerline{\begin{forest}
sm edges
[S
       [{NP[\snom]} [hvem;wer] ]
       [{S/NP[\snom]}
         [{NP[\lnom]} [det;\textsc{expl}] ]
         [{VP/NP[\snom]}
           [{\vbarSLASH NP[\snom]}
             [V [kende;kennt] ]
             [{NP[\snom]/NP[\snom]} [\trace] ] ]
           [{NP[\sacc]} [Kim;Kim ] ] ] ] ]
\end{forest}}
\caption{Analyse von \emph{hvem det knows Kim} `wer Kim kennt'}
%\caption{Analysis of interrogative clauses in Danish with subject extraction}\label{fig-interrogative-Danish-subject-extraction}
\end{figure}


\item Analysieren Sie den Satz in (\mex{1}). Verwenden Sie Dreiecke für die NP und die PP.
\ea
Es schwammen zwei Delphine neben dem Boot.
\z

\begin{figure}

\begin{forest}
sm edges
[S
  [{NP[\lnom]} [es]]
  [{S/NP[\lnom]}
     [V\sliste{ S//V } [schwammen]]
     [{S//\vSLASHnp[\lnom]}
       [{NP[\lnom]/NP[\lnom]}  [\trace]]
       [\vbarDSLv
         [{NP[\snom]} [zwei Delphine,roof]]
         [\vbarDSLv
           [PP [neben dem Boot,roof]]
           [V//V [\trace]]]]]]]
\end{forest}
\caption{Analyse von \emph{Es schwammen zwei Delphine neben dem Boot.}}
\end{figure}

\end{enumerate}



%      <!-- Local IspellDict: de_DE -->
